\documentclass[12pt,a4paper]{article}
\usepackage{cmap}
\usepackage[utf8]{inputenc}
\usepackage[T1,T2A]{fontenc}
\usepackage[english,russian]{babel}
\usepackage{amsmath,booktabs,array}
\usepackage{caption,float,qrcode}
\usepackage{geometry,hyperref,indentfirst}
\geometry{margin=25mm}
\hypersetup{unicode=true,colorlinks=true,urlcolor=blue,linkcolor=black}
\setlength{\parindent}{1.25cm}
\setlength{\parskip}{0pt}
\setlength{\emergencystretch}{2em}
\DeclareCaptionLabelSeparator{emdash}{\space---\space}
\captionsetup[figure]{justification=centering,singlelinecheck=false,labelsep=emdash}
\addto\captionsrussian{\renewcommand{\figurename}{Рис.}}
\linespread{0.96}
\begin{document}
\hypersetup{pageanchor=false}
\begin{titlepage}
\begin{center}
{\large МОСКОВСКИЙ АВИАЦИОННЫЙ ИНСТИТУТ\\
(НАЦИОНАЛЬНЫЙ ИССЛЕДОВАТЕЛЬСКИЙ УНИВЕРСИТЕТ)}
\vspace{18pt}
{\large Факультет информационных технологий и прикладной математики}

{\large Кафедра вычислительной математики и~программирования}
\vfill
{\large\textbf{Лабораторная работа №17 по курсу\\ «Информационный поиск»}}
\vspace{12mm}

{\large\textbf{Автотесты и модульное тестирование}}
\vfill
\begin{flushright}
\begin{tabular}{@{}r@{:\hspace{0.6em}}l@{}}
\textit{Студент} & К.\,А. Арусланов\\
\textit{Группа} & М8О-403Б-23\\
\textit{Преподаватель} & А.\,А. Кухтичев\\
\textit{Дата} & 08.11.2026\\
\textit{Оценка} &
\end{tabular}
\end{flushright}
\vfill
Москва, 2026
\end{center}
\end{titlepage}
\hypersetup{pageanchor=true}
\setcounter{page}{2}
\tableofcontents
\newpage

\section{Постановка задачи}
Подготовить воспроизводимый набор автоматических проверок для компонентов
учебной поисковой системы. Обосновать тестовый план, покрывающий основные
правила обработки документов, построения индекса и выполнения запросов.

\section{Метод решения}
Основные C++-компоненты проверяются GoogleTest. Тесты токенизатора
вызывают его непосредственно; индекс и поиск проверяются на небольших
синтетических корпусах с заранее известными результатами.
Тесты Python unittest проверяют правила робота, HTML-парсеры, экспорт,
веб-обвязку и обработку ошибок. Модульные тесты используют небольшие
входы с известным результатом. Интеграционные тесты запускаются с
отдельной временной MongoDB.

Общий запуск выполняет \verb|run_tests.py --lab all|: он пересобирает
C++-тесты в Docker, выполняет CTest по отдельным именованным случаям,
после чего запускает Python unittest в подробном режиме.
Скрипт \path{run_lab17.py} сохраняет полный стандартный вывод
и ошибки в \path{lab17/tests.txt}, запись о тестовом образе в
\path{lab17/build.txt}, журнал временной MongoDB в \path{lab17/mongo.txt}
и сводку с контрольными суммами в
\path{lab17/summary.json}. Если образа нет, он собирается из Dockerfile;
если есть, используется существующий. C++-программы при каждом прогоне
пересобираются из текущих исходников. Тестовая MongoDB запускается в
отдельном контейнере без постоянного тома и останавливается после прогона.

Для Variable Byte использован отдельный эталонный кодировщик, реализованный
в тесте по определению формата. На граничных значениях $0$, $127$, $128$,
$16383$, $16384$, $2^{32}$, $2^{63}$ и $2^{64}-1$ проверяется восстановление
числа из файла и из памяти. Декодер должен остановиться ровно после
текущего числа и не поглотить начало следующего. Каждый неполный префикс
десятибайтовой записи максимального числа должен быть отвергнут.

Для токенизатора проверено свойство объединения независимых полей.
Если между двумя корректными полями вставить разделитель, последовательности
терминов должны соединиться без изменений; позиции и байтовые границы
сдвигаются на известное число токенов и байтов. В проверке есть русские
слова, игровые имена, версия и C++/C\#.

Тесты Python проверяют правила обработки URL и метаданных. Канонизация
удаляет рекламные параметры, сохраняя значимую
пагинацию, и повторный вызов не меняет URL. Перенаправление StopGame
применяется только к числовому ID новости. Политика исключения отличает
страницу категории от вложенной публикации. В JSON-LD находится вложенная
публикация; повреждённый блок не скрывает следующий корректный.

\section{План тестирования}
План составлен по поведению компонентов и возможным отказам. Для каждой
группы указан независимый ожидаемый результат, а не только отсутствие
аварийного завершения.

\begin{center}
\small
\begin{tabular}{@{}p{29mm}p{61mm}p{52mm}@{}}
\toprule
Компонент & Проверяемое свойство & Ожидаемый результат\\
\midrule
Токенизатор & UTF-8, игровые имена, границы, соединение полей & Точные токены и байтовые смещения; ошибка на повреждённом UTF-8\\
Лемматизатор & Формы, неоднозначность, неизвестные слова & Все словарные варианты без выдуманных лемм\\
Индекс & Постинги, координаты, пустые документы и большие ID & Уникальные возрастающие списки, целостный формат\\
Поиск & AND/OR/NOT, фразы, порядок и соседство & Совпадение с таблицей истинности и независимой проверкой фразы\\
Сжатие & VB и контрольные точки & Точное восстановление; отказ на усечении и чужом индексе\\
Ранжирование & TF-IDF, нормировка, top-$k$ & Совпадение с независимой формулой и полной сортировкой\\
Робот & URL-категории, метаданные и ошибки & Статьи сохраняются, каталоги отсеиваются, неверные данные не принимаются\\
Интерфейсы & CLI и веб-форма & UTF-8 и корректный ответ после ошибочного запроса\\
\bottomrule
\end{tabular}
\end{center}

В наборе использованы фикстуры реальных
публикаций для токенизатора, небольшой словарь для морфологии,
таблицы истинности для булева поиска, собственное скользящее окно
для проверки фраз и прямой расчёт косинуса для TF-IDF.
Это позволяет находить ошибки в разных частях системы, не сравнивая
программу только с её же повторным запуском.

\section{Результаты}
\subsection{Запуск тестового плана}
Команда \verb|python run_lab17.py| выполнена из корня проекта.
Готовый Docker-образ использован повторно; исходники C++ собраны заново.
Тестовая MongoDB запущена отдельно и остановлена после проверок.

\begin{center}
\begin{tabular}{@{}lrrr@{}}
\toprule
Набор & Всего & Успешно & Пропущено\\
\midrule
CTest: GoogleTest и проверки CLI & 147 & 147 & 0\\
Python unittest & 291 & 291 & 0\\
\bottomrule
\end{tabular}
\end{center}

Все проверки завершились без ошибок. Шесть интеграционных тестов
работали с временной MongoDB; каждый использовал отдельную базу со случайным
именем. Журнал содержит 147 отдельных строк CTest с именами
проверок и результатами, а также имена и результаты 291 Python-теста.
Сводка сохраняет результаты и контрольные суммы файлов запуска.

Случаи выявляют определённые классы регрессий. Например, неверная
обработка границы 127/128 или усечённого VB-кода нарушит тест декодера;
потеря байтовой границы после соединения полей нарушит свойство
токенизатора. Ошибочное обобщение правила каталога PC Gamer приведёт
к отбрасыванию вложенной статьи, что покажет тест URL-политики.

Малые фикстуры не воспроизводят все сбои большого корпуса и внешних сайтов.
Однако они позволяют точно проверить поведение компонентов на известных входах.

\subsection{Воспроизводимость}
Исходники тестов находятся в \path{../tests/cpp/} и \path{../tests/python/}.
Команды и тестовый план приведены в \path{lab17_run.md}.
Журналы и сводка находятся в \path{lab17/}.
Для полной проверки используется \verb|python run_lab17.py|.

Исходный код:
\href{https://github.com/kirbasss/MAI-IR/tree/lab17}{github.com/kirbasss/MAI-IR/tree/lab17}.
\begin{figure}[H]
\centering
\begingroup
\hypersetup{urlcolor=black}
\qrcode[height=25mm]{https://github.com/kirbasss/MAI-IR/tree/lab17}
\endgroup
\caption{Ссылка на ветку лабораторной работы №17}
\end{figure}

\newpage
\section{Вывод}
В этой работе я собрал единый проверяемый план автотестов поисковой системы.
Тесты охватывают главные алгоритмы и пограничные случаи кодирования чисел,
токенизации, обработки URL и чтения метаданных. Интеграционные проверки
робота выполнены на отдельной временной MongoDB.
Полный прогон завершился успешно: 147 проверок CTest и 291 Python-тест.

Теперь результат можно повторить одной командой и проверить по именованному
журналу. Малый набор фикстур не отражает все состояния большого корпуса;
по обнаруженным сбоям следует добавлять новые случаи.
\end{document}
